
\documentclass[11pt,a4paper]{article}
\usepackage[utf8]{inputenc}
\usepackage[T1]{fontenc}
\usepackage[spanish]{babel}
\usepackage{geometry}
\usepackage{graphicx}
\usepackage{hyperref}
\usepackage{xcolor}
\usepackage{tabularx}
\geometry{margin=2cm}
\setlength{\parindent}{0pt}
\setlength{\parskip}{0.55em}
\sloppy
\definecolor{tcTeal}{HTML}{0E7C74}
\definecolor{tcDark}{HTML}{101817}
\definecolor{tcMuted}{HTML}{526A65}
\definecolor{tcWarnBg}{HTML}{FFF7E5}
\definecolor{tcWarnBorder}{HTML}{C99123}
\hypersetup{colorlinks=true, urlcolor=tcTeal, linkcolor=tcTeal}
\begin{document}

\begin{center}
{\huge\bfseries Informe AI Analyst: EURGBP.r H4 weavecount case}\\[0.25em]
{\large Reporte read-only para revision humana}
\end{center}

\vspace{0.2em}
\hrule height 0.7pt
\vspace{0.9em}

\begin{tabularx}{\textwidth}{>{\bfseries}p{3.2cm}X}
Paquete & EURGBP.r H4 W3 c6b5f358 \\
Setup & EURGBP.r H4 weavecount case \\
Modelo & gpt-5.5 \\
Thinking & medium \\
Macro & medium \\
Estado & reviewed \\
\end{tabularx}

\vspace{0.75em}
\fcolorbox{tcWarnBorder}{tcWarnBg}{%
\begin{minipage}{0.94\textwidth}
\textbf{Limite de uso.} Este informe no es una senal, no autoriza MT5 y no ejecuta ordenes. Sirve para priorizar revision humana y documentar el contexto.
\end{minipage}}

\vspace{0.4em}

\section*{Resumen}
Caso estructural H4 de EURGBP.r con lectura bajista candidata de calidad media. La estructura marca un posible tercer tramo a la baja desde 0.86562 hacia 0.86344, pero sigue en modo estudio: no es operativo, no autoriza ejecucion y requiere confirmacion adicional bajo 0.862 antes de considerar que el tramo gana validez estructural.

\section*{Lectura tecnica}
El recuento parte de un maximo en 0.86815 el 2026-05-29 12:00, desarrolla un primer tramo bajista hasta 0.862 el 2026-06-02 00:00, rebota hasta 0.86562 el 2026-06-04 16:00 y despues vuelve a ceder hasta 0.86344 al cierre H4 del 2026-06-05 16:00. Visualmente, el ultimo tramo bajista existe, pero aun queda por encima de la activacion marcada en 0.862, por lo que la lectura sigue siendo una hipotesis de continuidad y no una confirmacion completa.

\section*{Grafico de referencia}
\begin{center}
\includegraphics[width=0.96\textwidth]{\detokenize{ai_analyst_report_chart.png}}
\end{center}


\section*{Lectura de confluencias}
La direccion H4 y D1 aparece bajista, lo que favorece la coherencia temporal de una lectura estructural descendente en el marco principal; El cierre mas reciente, 0.86344, queda por debajo del punto de invalidacion 0.86562 y mantiene viva la hipotesis bajista desde el maximo de W2; El ultimo tramo bajista viene acompanado de incremento de actividad relativa en la vela final H4, con volumen de ticks 10052, superior al de varias velas previas del mismo tramo; y El RSI H4 en 37.06 y el RSI H1 en 43.61 no muestran sobreventa extrema, por lo que el movimiento bajista aun no aparece tecnicamente agotado por ese indicador.

\section*{Contradicciones y cautelas}
La activacion estructural en 0.862 no ha sido perforada al cierre del paquete; el precio termino en 0.86344, por encima del nivel que validaria mejor la continuidad bajista; La alineacion entre M15, H1 y H4 no esta sincronizada: M15 y H4 figuran bajistas, pero H1 figura alcista, lo que reduce limpieza direccional de corto plazo; El tramo W3? actual es corto frente al primer tramo bajista y aun no demuestra extension clara propia de un tercer movimiento impulsivo; y La volatilidad H4 esta por debajo de su mediana reciente, con ATR relativo H4 en 0.10 frente a mediana 0.17; esto sugiere menor expansion del rango en el marco principal.

\section*{Riesgos principales}
La lectura es exclusivamente de estudio estructural y no debe tratarse como instruccion operativa; El precio se encuentra entre invalidacion 0.86562 y activacion 0.862, una zona intermedia donde el sesgo visual puede cambiar con pocas velas H4; La calidad media obliga a revisar manualmente la proporcionalidad de W1, W2 y W3?, especialmente porque el tramo actual todavia no confirma ruptura del minimo previo; y El par EURGBP puede reaccionar de forma brusca a repricing relativo entre BCE y Banco de Inglaterra, por lo que el contexto macro aumenta la cautela de revision.

\section*{Contexto macro y noticias}
El contexto macro reciente para EURGBP esta dominado por una comparacion de bancos centrales. En la zona euro, Eurostat publico que la inflacion anual estimada subio a 3.2\% en mayo de 2026, con energia y servicios como componentes relevantes. Esa lectura sostiene la sensibilidad del euro a expectativas de endurecimiento del BCE, especialmente porque Reuters habia informado que el mercado seguia pendiente de una posible subida de tipos en junio.

Para EURGBP, esto introduce riesgo de movimientos por repricing de tipos, no solo por estructura tecnica. En Reino Unido, el Banco de Inglaterra mantiene el tipo bancario en 3.75\% y su proxima decision esta prevista para el 18 de junio de 2026, mientras la inflacion oficial vigente publicada por el BoE aparece por encima del objetivo. El cruce puede quedar atrapado entre dos narrativas: euro apoyado por presion inflacionista y posible BCE mas restrictivo, frente a libra condicionada por inflacion britanica, expectativas de recortes y sensibilidad a energia/riesgo geopolitico. Por tanto, el riesgo macro para revisar el setup es medio: no invalida la lectura tecnica, pero aconseja cautela porque cualquier sorpresa de politica monetaria o inflacion puede desplazar el par fuera de los niveles estructurales.

Desde una lectura financiera, este contexto no cambia por si solo el setup tecnico, pero si condiciona la prudencia de la revision: aumenta el peso de comprobar calendario, liquidez y confirmacion posterior antes de extraer conclusiones. Con riesgo macro marcado como medio, la lectura debe tratarse como sensible a evento y no como una validacion aislada del grafico.

\section*{Conclusion operativa y financiera}
Conclusion operativa y financiera: EURGBP.r queda como weavecount case de revision, no como instruccion de mercado. La prioridad de revision indicada por el analisis es 3/5 y el riesgo macro figura como medio. La interpretacion conjunta debe apoyarse en tres filtros humanos: que el precio confirme el comportamiento descrito en el grafico, que las contradicciones de tendencia o estructura no invaliden el caso, y que no exista un evento macro cercano capaz de distorsionar el movimiento. Sin esas comprobaciones, el informe solo justifica seguimiento y documentacion, no ejecucion.

\section*{Comprobaciones humanas}
\begin{itemize}
  \item Comprobar si hay cierre H4 limpio por debajo de 0.862 o rechazo claro antes de ese nivel.
  \item Revisar si un cierre por encima de 0.86562 invalida la lectura bajista candidata.
  \item Comparar el recuento H4 con H1 para ver si el tramo bajista actual tiene subdivisiones coherentes o si es solo correccion interna.
  \item Verificar proximidad de datos de inflacion, decision del BCE o decision del Banco de Inglaterra antes de cualquier uso del caso en paneles de seguimiento.
\end{itemize}

\section*{Fuentes}
\textbf{Fuentes locales}\par
\begin{tabularx}{\textwidth}{p{4.2cm}X}
setup\_context.json & Datos estructurados del setup \\
market\_context.json & Contexto agregado de mercado \\
chart\_layers.csv & Capas dibujadas en el grafico \\
ohlc\_window.csv & Ventana de velas OHLC \\
source\_manifest.json & Trazabilidad del paquete \\
chart.png & Imagen renderizada del grafico \\
\end{tabularx}

\vspace{0.45em}
\textbf{Enlaces externos}\par
\begin{tabularx}{\textwidth}{p{7.4cm}X}
\href{\detokenize{https://ec.europa.eu/eurostat/en/web/products-euro-indicators/w/2-02062026-ap}}{Eurostat, flash estimate May 2026} & Fuente externa consultada para contexto macro/noticias. \\
\href{\detokenize{https://www.bankofengland.co.uk/monetary-policy}}{Bank of England monetary policy page} & Fuente externa consultada para contexto macro/noticias. \\
\href{\detokenize{https://www.bankofengland.co.uk/monetary-policy-summary-and-minutes/2026/april-2026}}{Bank of England April 2026 Monetary Policy Summary} & Fuente externa consultada para contexto macro/noticias. \\
\href{\detokenize{https://www.ons.gov.uk/releases/consumerpriceinflationukapril2026}}{Office for National Statistics, UK CPI April 2026 release} & Fuente externa consultada para contexto macro/noticias. \\
\href{\detokenize{https://www.marketscreener.com/news/ecb-june-rate-hike-case-is-nearly-sealed-but-july-is-fully-open-sources-say-ce7f5ad9da8eff20}}{Reuters via MarketScreener, ECB June hike expectations} & Fuente externa consultada para contexto macro/noticias. \\
\end{tabularx}

\end{document}
